\section{Discussion and Conclusions}
\label{sec:conclusions}
%% ============================================================

\subsection{The Story in Brief}

We started with the EGPT design: gradient descent on an implicit energy function,
realised as a $V{=}K$ attention constraint plus learned projection.
The original design is compute-inefficient (4.4$\times$ FLOPs overhead) and trails GPT
on perplexity despite tying on accuracy.

Mechanistic analysis reveals that EGPT learns a balanced Helmholtz decomposition
(curl $\approx$ gradient) in its projections, and that EGPT attention outputs
are more directionally aligned across layers than in GPT --- consistent with the
shared-energy hypothesis.
However, this functional similarity does not imply structural redundancy:
weight merging is lossy even for the most similar adjacent layer pairs.

Mixed heads (energy $V{=}K$ + GPT $V{=}W_V x$) close the perplexity gap with GPT
while matching accuracy, but a GSM8K regression reveals the $V{=}K$ output rule
introduces a bias against arithmetic reasoning.

The key insight is that the energy-head output under the $V{=}K$ rule is not the
true energy gradient: it lies on the key manifold rather than pointing in the direction
of steepest descent in hidden space.
Replacing it with the true gradient ($W_Q^\top \cdot \mathrm{softmax}(\cdot) K$)
--- EGrad --- simultaneously improves accuracy, perplexity, and GSM8K, strictly
dominating GPT at 144M scale.
At 400M scale, true EGrad (V31) achieves 50.7\% avg / 27.9 PPL vs.\ GPT's 48.6\% / 29.8 PPL
at identical FLOPs --- the first consistent energy-model improvement over GPT.

\subsection{What Makes EGrad Work?}

The $W_Q^\top$ factor in the EGrad output plays two roles:
(1) it maps from head space ($d_h$) to hidden space ($D$), allowing the energy-head
signal to interact with the full-dimensional residual stream;
(2) it introduces a rotation/scaling that is learned alongside $W_Q$, effectively
giving the energy head a free ``output projection'' at no parameter cost.
EDesc achieves similar improvements by the subtractive update alone (no $W_Q^\top$),
suggesting that the directionality of the update matters as much as the exact gradient form.

\subsection{Open Questions}

\begin{enumerate}
  \item \textbf{Test-time scaling.}  EGrad models are trained with a fixed number of
    recurrence steps.  Can extra steps at test time (without retraining) improve
    accuracy on harder tasks?  The shared-energy structure motivates iterative refinement.
  \item \textbf{Recurrence schedule.}  V25/V26 will test whether concentrating iterations
    in later layers is better than uniform distribution.  If so, a learned per-layer
    schedule may be optimal.
  \item \textbf{Larger scale.}  The EGrad advantage at 400M is $\sim$1.3 pp accuracy
    and $\sim$2 PPL.  Does this gap grow or shrink at 3B+ scale?
  \item \textbf{RL-steered inference.}  The EGrad energy gradient provides a natural
    reward signal for RL-based test-time search: REINFORCE over recurrence schedules.
\end{enumerate}

\subsection{Summary of Findings}

\begin{enumerate}
  \item EGPT's $4.4\times$ FLOPs overhead is not justified by accuracy gains.
  \item Trained EGPT projections learn balanced Helmholtz decomposition ($\|S\|_F \approx \|A\|_F$);
    enforcing it architecturally hurts training.
  \item EGPT attention layers share a common gradient direction (cosine 0.176 vs GPT 0.119),
    consistent with a shared energy function; FFN layers do not.
  \item High layer-similarity does not imply weight-mergeability; functional similarity is
    decoupled from structural similarity.
  \item Mixed heads (6E + 6G per block) match GPT accuracy at fewer FLOPs but regress on GSM8K.
  \item Replacing the $V{=}K$ energy-head output with the true energy gradient (EGrad)
    strictly dominates GPT on all metrics at 144M scale.
  \item At 400M / 503 MFLOPs/tok, true EGrad (V31) achieves 50.7\% avg / 27.9 PPL,
    outperforming matched GPT (V9: 48.6\% / 29.8 PPL) by $+2.1$ pp accuracy and $-1.9$ PPL,
    and narrowly exceeding Mixed V19 (49.9\% / 27.8 PPL) --- the first consistent energy-model win over GPT.
  \item Weight-shared recurrent EGrad (V23/V24) bridges the gap but costs $\sim$4 PPL
    vs.\ non-recurrent at matched compute.
\end{enumerate}

